\documentclass{ltxdoc}
\ProvidesFile
    {nyu-logos.tex}
    [2023/05/29 v0.11c Logos in the Gotham bundle]
\GetFileInfo{nyu-logos.tex}
\title{Logos in the Gotham bundle}
\author{Matthew Leingang\thanks{leingang@nyu.edu}} 
\date{\fileversion, Released \filedate}

\usepackage{titlesec}
\usepackage{titling}
\usepackage{fontspec}
\newfontfamily\headingfont[]{Gotham Medium}
\titleformat*{\section}{\Large\headingfont}
\titleformat*{\subsection}{\large\headingfont}
\titleformat*{\subsubsection}{\headingfont}
\renewcommand{\maketitlehooka}{\headingfont}
\setmainfont{Frank Ruhl Libre}
\setsansfont{Gotham Book}
% \setmonofont{Decima Mono Cyr}

\usepackage{bxtexlogo}\bxtexlogoimport{*}
\usepackage[letterpaper]{geometry}
\usepackage{adjustbox} % allow normally overfull hboxes to be centered.
\usepackage{url}% \url

\usepackage{graphicx}
\usepackage{tikz}
\usetikzlibrary{patterns}

\newsavebox{\logobox}
\newcommand\showwhite[1]{
    \savebox{\logobox}{#1}
    \begin{tikzpicture}[baseline=0.1\ht\logobox]
        \fill[gray] (0,0) rectangle (1.2\wd\logobox,1.2\ht\logobox);
        \node[anchor=south west,inner sep=0pt,outer sep=0pt] at (0.1\wd\logobox,0.1\ht\logobox) {
            \usebox{\logobox}
        };
    \end{tikzpicture}
}

\begin{document}
\maketitle

\section{Introduction}

This document describes image files in the \texttt{logos} module within the 
Gotham bundle.  The Gotham bundle is for \LaTeX{} users to produce documents 
conforming to NYU's visual identity.  

\section{Installation}

The bundle itself can be downloaded from Github.  Clone the repository,
or unzip a release package, and within the directory execute \texttt{l3build install --full}.
This will install all the modules within the bundle, including this one.

There is no \LaTeX{} package installed by this module.  Instead, the image files 
are copied to the user's TEXMFHOME tree, so that they can be included anywhere 
by use of the \verb|\includegraphics| command of the \texttt{graphicx} bundle.

\section{Usage}

Make sure your document loads the \texttt{graphicx} package by having
\verb|\usepackage{graphicx}| in the preamble.  Then at any part of the document,
put \verb|\includegraphics|\marg{logo}, where \marg{logo} is any of the logos
named here.

You don't need a file extension.  \LaTeX{} will look for one compatible with the
engine.  Which one is selected first is determined by the driver files.  See
\cite{tse-graphics-order} for the details.

Most of the time you want PDF if you're producing a PDF.  PDF is a vector format,
meaning it scales to arbitrary size perfectly.  With bit-mapped formats like PNG,
you will get jagged edges at large scale.  JPG is more for image with thousands
of colors like photos, than it is for two-color images like logos.

You can always check the log file to see which file is actually included, and force
a different extension if it's not the one you want.


\section{The Logos}

All available logos are displayed in this section, at their natural size.  Images 
that are all in white are shown over a gray background, but the background is not 
part of the image.

The university's visual identity guide \cite{nyu-downloads} (see also
\cite{nyu-logo-basics}) offers this advice on which logo “shape” to use when:

\begin{itemize}
    \item The short logo is used primarily for internal communications, like for 
    lecture slides or a departmental presentation.
    \item The long logo is used when extra formality is needed, or “for outside
    audiences where the full name adds additional information.”  A conference presentation,
    perhaps.
    \item The stacked logo for small spaces, like the corner of a slide.
\end{itemize}

I found EPS logos in a different shade of purple than the rest.  I corrected these
EPS files to use the official purple, and those are the sources for the PDF logos. 
Since I used Inkscape to do this, and Inkscape's native format is SVG, there are 
some SVG logo files in the module as well.  

\subsection{NYU “Flagship” logos}

See Figure~\ref{fig-nyu-logos}.

\begin{figure}
    \centering
    \begin{adjustbox}{center}
    \begin{tabular}{ccc}
        \includegraphics{nyu_long_black} &
        \includegraphics{nyu_long_color} &
        \showwhite{\includegraphics{nyu_long_white}}
        \\
        \texttt{nyu\_long\_black} &
        \texttt{nyu\_long\_color} &
        \texttt{nyu\_long\_white} \\[2em]
        \includegraphics{nyu_short_black} &
        \includegraphics{nyu_short_color} &
        \showwhite{\includegraphics{nyu_short_white}} \\
        \texttt{nyu\_short\_black} &
        \texttt{nyu\_short\_color} &
        \texttt{nyu\_short\_white} \\[2em]
        \includegraphics{nyu_stacked_black} &
        \includegraphics{nyu_stacked_color} &
        \showwhite{\includegraphics{nyu_stacked_white}} \\
        \texttt{nyu\_stacked\_black} &
        \texttt{nyu\_stacked\_color} &
        \texttt{nyu\_stacked\_white}
    \end{tabular}
    \end{adjustbox}
    \caption{%
        The main NYU logos
    }
    \label{fig-nyu-logos}
\end{figure}


\subsection{Arts and Science logos}


\subsection{Courant logos}

See Figure~\ref{fig-cims-logos}.

\begin{figure}
\begin{adjustbox}{center}
    \begin{tabular}{cc}
        \includegraphics{courant_short_black} &
        \showwhite{\includegraphics[width=2.8194444444in]{courant_short_white.png}} \\
        \texttt{courant\_short\_black} &
        \texttt{courant\_short\_white.png} 
    \end{tabular}
    \end{adjustbox}
    \caption{%
        NYU Courant logos.
        At the moment \texttt{courant\_short\_white.pdf} is broken, so we only
        include the PNG version.
    }
    \label{fig-cims-logos}
\end{figure}

\subsection{Tandon School of Engineering logos}

\section{License and Legalishness}

The NYU Logo is a trademark of New York University.  You can't use the logo to
create an impression that your products are affiliated or endorsed by the University,
except in the cases that they are.

The logos were downloaded from publicly available web pages (\cite{nyu-downloads}).
% Also https://as.nyu.edu/communications/visual-identity-downloads.html
Their aggregation in this module does not constitution of a claim of origination
from the module authors.

Gotham is a common nickname for New York City.  Gotham is also the name of a typeface
created by Hoefler\&Co, which happens to be the official logo font of NYU.
But this bundle is not affiliated with Hoefler\&Co.

\bibliographystyle{plain}
\bibliography{nyu-logos}


\end{document}